\subsection{Przegląd technik głębokiego uczenia i wyzwania związane ze zbiorami danych w analizie twarzy przysłoniętych}

\paragraph{}Kompleksowe ujęcie problemu analizy twarzy z okluzją wymaga uwzględnienia nie tylko architektury sieci, ale także specyfiki danych treningowych i systematyki stosowanych metod. W najnowszym artykule przeglądowym z 2025 roku, Thaher i in. [5] wskazują, że okluzja stanowi jedno z najważniejszych otwartych wyzwań dla systemów biometrycznych i nadzoru wizyjnego, drastycznie wpływając na skuteczność algorytmów detekcji i rozpoznawania.

Autorzy dokonują systematyzacji podejść badawczych, wyróżniając metody oparte na głębokim uczeniu jako dominujący standard. Wskazują przy tym na dwa główne nurty radzenia sobie z problemem:

1. Metody oparte na danych (Data-centric): Polegające na augmentacji zbiorów treningowych poprzez syntetyczne nakładanie masek lub innych obiektów przysłaniających na twarze, co pozwala modelom "nauczyć się" okluzji.

2. Metody oparte na modelu (Model-centric): Obejmujące projektowanie specyficznych architektur, takich jak sieci wykorzystujące mechanizmy uwagi (np. wspomniany wcześniej CBAM) czy piramidy cech (FPN), które pozwalają na ekstrakcję informacji z widocznych fragmentów twarzy mimo zakłóceń.

\paragraph{}Istotnym aspektem poruszonym w przeglądzie jest rola specjalistycznych zbiorów danych. Thaher i in. podkreślają, że standardowe bazy twarzy są niewystarczające do trenowania odpornych systemów. Wymieniają kluczowe zbiory dedykowane problemowi okluzji, takie jak MAFA (Masked Faces), które zawierają twarze z naturalnymi przesłonięciami, a nie tylko syntetycznymi. Autorzy zwracają uwagę, że sukces w implementacji modeli odpornych na okluzję zależy w dużej mierze od jakości i różnorodności tych danych, a także od zdolności algorytmów do radzenia sobie z ekstremalnymi warunkami, takimi jak połączenie okluzji z trudnym oświetleniem czy nietypową pozą [5].

Wnioski z tego przeglądu uzasadniają konieczność testowania implementowanych rozwiązań nie tylko na ogólnych zbiorach walidacyjnych, ale także na dedykowanych podzbiorach zawierających "twardą" okluzję, aby rzetelnie ocenić ich przydatność w rzeczywistych zastosowaniach.

Pomimo intensywnych badań nad rozpoznawaniem twarzy w warunkach okluzji, literatura wciąż w ograniczonym stopniu analizuje wpływ lekkich, modułowych mechanizmów atencji na odporność nowoczesnych architektur opartych na funkcjach strat z marginesem kątowym. W szczególności brakuje jednoznacznych analiz dotyczących umiejscowienia takich mechanizmów w strukturze sieci oraz ich wpływu na skuteczność identyfikacji w warunkach syntetycznej i rzeczywistej okluzji.